\section{Linear Algebra}

Here, we just include many problems that will be useful. Some exercises
emphasize how some wishful statements do not hold.

\begin{exercise}
    Show that $\{1, (x-1), (x-1)(x-2)\}$ is a basis of $\P_3[x]$ and that
    $W = \{p(x) \in \P_3[x] \colon p(1) = 0\}$ is a subspace of $\P_3[x]$. Find $\dim W$.
    Here, $\P_3[x]$ denotes the space of real polynomials of degree strictly less than $3$.
\end{exercise}
\begin{solution}
    The fact that this is a basis is clear. We find $\dim W$ using rank-nulity.
    The image has dimension $1$, since two of the basis elements vanish. 
    Thus, $\dim W=2$.
\end{solution}

\begin{exercise}
    Let $\alpha_1,\alpha_2,\alpha_3,$ and $\alpha_4$ be linearly independent. 
    Answer true or false.
    \begin{enumerate}
        \item $\alpha_1+\alpha_2, \alpha_2+\alpha_3, \alpha_3+\alpha_4, \alpha_4+\alpha_1$ are linearly independent.
        \item $\alpha_1-\alpha_2, \alpha_2-\alpha_3, \alpha_3-\alpha_4, \alpha_4-\alpha_1$ are linearly independent.
        \item $\alpha_1+\alpha_2, \alpha_2+\alpha_3, \alpha_3+\alpha_4, \alpha_4-\alpha_1$ are linearly independent.
        \item $\alpha_1+\alpha_2, \alpha_2+\alpha_3, \alpha_3-\alpha_4, \alpha_4-\alpha_1$ are linearly independent.
    \end{enumerate}
\end{exercise}
\begin{solution}
    1. False. We can make an easy relation among these elements.\\

    2. False. Just sum them all. \\

    3. True. Just expand. \\

    4. False. Sum of last two is equal to difference of second and first.
\end{solution}

\begin{exercise}
Let $t \in \R$. Discuss the linear independence of the vectors over $\R$:
\[
\alpha_1=(1,1,0), \qquad
\alpha_2=(1,3,-1), \qquad
\alpha_3=(5,3,t).
\]
\end{exercise}
\begin{solution}
We have the relations:
\[
a + b + 5c = 0, \qquad a + 3b + 3c = 0, \qquad -b + tc = 0
\]
The first two equations provide that $2b - 2c = 0$, or $b = c$. Thus,
so long as $t \neq 1$, these are LI. 
More systematically, use the determinant condition.
\end{solution}

\begin{exercise}
Let $W_1$ and $W_2$ be nontrivial subspaces of a vector space $V$. The sum
$W_1+W_2$ is called a \emph{direct sum}, denoted as $W_1\oplus W_2$, if every element
$\alpha\in W_1+W_2$ can be uniquely written as
\[
\alpha=w_1+w_2,
\]
where $w_1\in W_1$ and $w_2\in W_2$. Show that the following statements are equivalent:
\begin{enumerate}
\item $W_1+W_2$ is a direct sum.    
\item If $w_1+w_2=0$ and $w_1\in W_1$, $w_2\in W_2$, then $w_1 = w_2 = 0$.
\item $W_1\cap W_2=\{0\}$.
\item $\dim(W_1+W_2)=\dim W_1+\dim W_2$.
\end{enumerate}

How can the direct sum be extended to more than two subspaces?
\end{exercise}
\begin{solution}
Standard.
\end{solution}

\begin{exercise}
Let $A,B,C$ be matrices. Prove if $r(A)$ denotes the rank of $A$, then 
prove the following:
\begin{enumerate}
    \item $r(A+B)\leq r(A)+r(B)$.
    \item $r(AB)\leq \min\{r(A),r(B)\}$.
    \item $r(AB)+r(BC)-r(B)\leq r(ABC)$.
\end{enumerate}
\end{exercise}
\begin{solution}

\end{solution}

\begin{exercise}
Let $A \in M_n(\R)$.
\begin{enumerate}
    \item Show that if $A^t=-A$ and $n$ is odd, then $\det(A) = 0$.
    \item Show that if $A^2 + I = 0$, then $n$ is even. What is $A \in M_n(\C)$ instead?
\end{enumerate}
\end{exercise}
\begin{solution}
    1. We have $\det(A) = -\det(A)$, becaause $n$ is odd, so $\det(A) = 0$. \\

    2. We have $\det(A)^2 = (-1)^n$. Therefore, $n$ must be even. If $A$ is complex,
    then $n$ can be even, because $\det(A) \in \C$ now.
\end{solution}

\begin{exercise}
Let $A, B \in M_n(\R)$. Show that
\[
\det 
\begin{pmatrix}
A & B\\
-B & A
\end{pmatrix} \geq 0.
\]
State the analogue for complex matrices.
\end{exercise}
\begin{solution}
$\det(A)^2 + \det(B)^2$ is a sum of nonnegative reals, so is nonnegative. If we 
want a similar statement for complex matrices, then we just need to conjugate the bottom two:
\[
\det 
\begin{pmatrix}
A & B\\
\overline{-B} & \overline{A}
\end{pmatrix} \geq 0
\]
\end{solution}

\begin{exercise}
Let $A, B \in M_n(\C)$. Show that
\[
\det
\begin{pmatrix}
A & B\\
B & A
\end{pmatrix}
=
\det(A + B) \det(A - B)
\]
\end{exercise}
\begin{solution}
    Just use the fact that $\det(A^2 - B^2) = \det((A - B)(A + B)) = \det(A - B)\det(A+B)$.
\end{solution}

\begin{exercise}
For matrices of appropriate sizes, answer true or false:
\begin{enumerate}
\item If $A^2 = B^2$, then $A = B$ or $A = -B$.
\item If $r(A) = r(B)$, then $r(A^2) = r(B^2)$.
\item $r(A + kB) \leq r(A) + kr(B)$, where $k$ is a positive scalar.
\item $r(A - B) \leq r(A) - r(B)$.
\item If $r(AB) = 0$, then $r(BA) = 0$.
\item If $r(AB) = 0$, then $r(A) = 0$ or $r(B) = 0$.
\end{enumerate}
\end{exercise}
\begin{solution}
1. False. Take $A = \begin{pmatrix}
    1 & 0 \\ 0 & -1
\end{pmatrix}, B =I$. \\

2. Take $A = \begin{pmatrix}
    0 & 1 \\ 0 & 0
\end{pmatrix}, B = \begin{pmatrix}
    1 & 0 \\ 0 & 0
\end{pmatrix}$. \\

3. False, take $A - 0$, $B \neq 0$, and $k = \frac{1}{2}$. The correct statement instead
\[
r(A + kB) \leq r(A) + kr(B).
\]

4. Take $A = 0$, $B \neq 0$. Then, $r(-B)$ is positive, while $-r(B)$ is negative. \\

5. Take $A = \begin{pmatrix}
    1 & 0 \\ 0 & 0
\end{pmatrix}$, $B = \begin{pmatrix}
    0 & 0 \\ 1 & 0
\end{pmatrix}$. Then, $AB = 0$, while $BA = B \neq 0$. \\

6. Use the same example as 2.
\end{solution}

\begin{proposition}
Show that $AB = I$ implies $BA = I$, where $A$, $B$ are matrices over a field.
\end{proposition}
\begin{solution}
This relies on some knowledge of algebra; particularly, about left or right Artinian rings. 
Note that the ring of matrices (over a field) form a finite dimensional algbera, and so
they form an Artinian ring. We recall that left or right Artinian rings are rings where
the DCC holds for left or right ideals. Let us  use the left version. 
The claim we prove is that $ab = 1$ implies $ba = 1$, for any elements $a,b$ 
(which satisfies the desired condition). We let the Artinian ring be called $R$. We note that
$R \subseteq bR \subseteq b^2R \subseteq \dots$ is a descending chain of left 
ideals, and so $b^k R = b^{k + 1}R$ for some positive integer $k$.
In other words, $b^k = b^{k + 1}c$, for some $c \in R$. Then, if we multiply by $a^k$ 
on the left, we get that $bc = 1$. Then, $ba = babc = bc = 1$.
\end{solution}

\begin{exercise}
Let $v_1, v_2, v_3, v_4$ be vectors in $\mathbb{R}^2$, and consider the following six dot products:
\[
v_1 \cdot v_2,\qquad
v_1\cdot v_3,\qquad
v_1\cdot v_4,\qquad
v_2\cdot v_3,\qquad
v_2\cdot v_4,\qquad
v_3\cdot v_4.
\]
Which of the following statements about the six dot products could be true?
\begin{enumerate}
    \item All six are negative.
    \item Two are negative, and four are equal to $0$.
    \item Two are positive, and four are equal to $0$.
\end{enumerate}
\end{exercise}
\begin{solution}
Testing with some random vector might be good, but takes some clever guesses
which I don't have the time to be searching for. The key is the dot product 
identity, which measures the angle depending on the sign. If positive, then acute,
if $0$, then perpendicular, if negative, then obtuse. Given four vectors 
in $\R^2$, then all six are negative if all vectors are obtuse with respect to each 
other. This is not possible. The remaining cases are possible.
In the case 2, we situate them all perpendicularly to each other. In the case 3,
we let $v_1 = v_2$, $v_3 = v_4$, and make these two vectors perpendicular to each other.
\end{solution}